% Acta N°8 - Consejo de Carrera CEIC - Jueves 8 de octubre de 2026
% Redactada a partir de los apuntes tomados en sesión (incompletos); lo marcado con \dato{...} debe verificarse o completarse.
%
% Compilación:
%   Versión final:     pdflatex Acta_CEIC_Jueves_8_de_Octubre.tex   (dos veces)
%   Versión revisión:  pdflatex -jobname=Acta_CEIC_Jueves_8_de_Octubre_REVISION "\def\revision{}% Acta N°8 - Consejo de Carrera CEIC - Jueves 8 de octubre de 2026
% Redactada a partir de los apuntes tomados en sesión (incompletos); lo marcado con \dato{...} debe verificarse o completarse.
%
% Compilación:
%   Versión final:     pdflatex Acta_CEIC_Jueves_8_de_Octubre.tex   (dos veces)
%   Versión revisión:  pdflatex -jobname=Acta_CEIC_Jueves_8_de_Octubre_REVISION "\def\revision{}% Acta N°8 - Consejo de Carrera CEIC - Jueves 8 de octubre de 2026
% Redactada a partir de los apuntes tomados en sesión (incompletos); lo marcado con \dato{...} debe verificarse o completarse.
%
% Compilación:
%   Versión final:     pdflatex Acta_CEIC_Jueves_8_de_Octubre.tex   (dos veces)
%   Versión revisión:  pdflatex -jobname=Acta_CEIC_Jueves_8_de_Octubre_REVISION "\def\revision{}% Acta N°8 - Consejo de Carrera CEIC - Jueves 8 de octubre de 2026
% Redactada a partir de los apuntes tomados en sesión (incompletos); lo marcado con \dato{...} debe verificarse o completarse.
%
% Compilación:
%   Versión final:     pdflatex Acta_CEIC_Jueves_8_de_Octubre.tex   (dos veces)
%   Versión revisión:  pdflatex -jobname=Acta_CEIC_Jueves_8_de_Octubre_REVISION "\def\revision{}\input{Acta_CEIC_Jueves_8_de_Octubre.tex}"
%
% En la versión de revisión, todo lo marcado con \dato{...} aparece en color:
% son datos que deben verificarse o completarse antes de publicar el acta.

\documentclass[12pt,letterpaper]{article}

\usepackage[T1]{fontenc}
\usepackage[utf8]{inputenc}
\usepackage[spanish,es-nolists,es-nodecimaldot]{babel}
\usepackage[letterpaper,left=1in,right=1in,top=1.35in,bottom=1.1in,headheight=48pt,headsep=24pt,footskip=36pt]{geometry}
\usepackage{graphicx}
\usepackage{amssymb}
\usepackage{xcolor}
\usepackage{fancyhdr}
\usepackage{lastpage}
\usepackage{titlesec}
\usepackage{enumitem}
\usepackage{textcomp}

% ---------------------------------------------------------------------------
% Datos a verificar
% ---------------------------------------------------------------------------
\newif\ifrevision
\ifdefined\revision \revisiontrue \else \revisionfalse \fi
\newcommand{\dato}[1]{\ifrevision\textcolor[HTML]{C0392B}{#1}\else#1\fi}

% ---------------------------------------------------------------------------
% Datos del acta
% ---------------------------------------------------------------------------
\newcommand{\numeroacta}{8}
\newcommand{\fechaacta}{08-10-2026}

% ---------------------------------------------------------------------------
% Formato
% ---------------------------------------------------------------------------
\setlength{\parindent}{1.5em}
\setlength{\parskip}{6pt}

\pagestyle{fancy}
\fancyhf{}
\fancyhead[L]{\raisebox{-0.35\height}{\includegraphics[height=1.15cm]{logo_ceic.png}}}
\fancyhead[R]{Acta N\textdegree\numeroacta\\\fechaacta}
\fancyfoot[C]{Página \thepage\ de \pageref{LastPage}}
\renewcommand{\headrulewidth}{0.4pt}
\renewcommand{\footrulewidth}{0.4pt}

\titleformat{\section}{\normalfont\Large\bfseries}{\thesection.}{0.35em}{}
\titleformat{\subsection}{\normalfont\large\bfseries}{\thesubsection.}{0.35em}{}
\titlespacing*{\section}{0pt}{3.0ex plus 1ex minus .2ex}{1.6ex plus .2ex}
\titlespacing*{\subsection}{0pt}{2.4ex plus 1ex minus .2ex}{1.2ex plus .2ex}

\renewcommand{\labelitemi}{\tiny$\blacksquare$}
\setlist[itemize]{leftmargin=2.2em,itemsep=2pt}
\setlist[enumerate]{leftmargin=2.2em,itemsep=2pt}

\begin{document}

% ---------------------------------------------------------------------------
% Encabezado
% ---------------------------------------------------------------------------
\begin{center}
  {\Large\bfseries Centro de Estudiantes de Ingeniería Comercial}\\
  Universidad de Chile, Facultad de Economía y Negocios

  \vspace{1.2em}
  \textbf{Acta Consejo}
\end{center}

\vspace{0.8em}

\textbf{Asistentes:} Alan Mc Coll (Mesa), Diego Pakomio (Mesa extendida), Anabelee Araneda (DM), Arturo Mora (DM), Juan Budge (DM), Martín Farías (DPC), Florencia Montiglio (DPC), Sofía Cisternas (DA), Nicolás Sanzana (DE), Sofía Allende (CT).

\textbf{Hora de inicio:} 14:00 hrs

\textbf{Hora de término:} 15:20 hrs

\vspace{1em}

\textbf{Tabla:}

\begin{enumerate}[itemsep=6pt]
  \item Finanzas transparentes
  \item Solemnes
  \item Jaranazos
  \item Semana organizaciones
  \item JOE
  \item Reforma estatutos
  \item Varios
\end{enumerate}

% ---------------------------------------------------------------------------
\section{Finanzas transparentes}

\subsection{Balance presentado}

Se presenta al Consejo el balance correspondiente a los meses de Agosto y Septiembre, junto con el detalle de las principales fuentes de ingreso y categorías de gasto registradas en el libro diario de cada mes.

\begin{itemize}
  \item \textbf{Saldo del mes de Agosto (Ingresos $-$ Gastos):} \$161.503
  \item \textbf{Saldo del mes de Septiembre (Ingresos $-$ Gastos):} \$-682.811
\end{itemize}

Lo que arroja un saldo actual de caja chica de: \$9.820.820

\subsection{Detalle del periodo}

En Agosto, los ingresos ascendieron a \$261.840 y provinieron principalmente de las entradas de Cupido FEN (\$129.850), de la repartición de ganancias del Jaranazo (\$101.990) y de la participación del CEIC en la charla de diversidad realizada junto a Fundación Todo Mejora (\$30.000). Los gastos del mes alcanzaron los \$100.337 y correspondieron a la realización de Cupido FEN y del Jaranazo, a las libretas y stickers de la SUE para la feria de organizaciones y a gastos menores de funcionamiento.

En Septiembre, los ingresos fueron de \$221.550, compuestos por las entradas SIPEE (\$147.000) y por la devolución de una transferencia realizada por error (\$74.550). Los egresos sumaron \$904.361, entre los cuales destacan los gastos de la FENDA (guardias y DJ, entre otros), las dos jornadas de Tiny Desk, las libretas y stickers de la SUE, Cupido FEN y los desayunos realizados para deportistas y funcionarios.

\textbf{Nota:} Cinco de las entradas SIPEE se pagaron a \$7.000 en lugar de \$8.000 por un error de formulario; los \$1.000 restantes de cada una aparecerán en el balance de Octubre. Se recuerda que cualquier estudiante puede solicitar el detalle y los comprobantes de los balances presentados a través del formulario del Portal de Transparencia del CEIC.

% ---------------------------------------------------------------------------
\section{Solemnes}

\subsection{Antecedentes}

Se abre la discusión sobre la distribución de las fechas de las solemnes del semestre y su impacto en las distintas menciones y generaciones, así como sobre la coordinación general de las evaluaciones con las clases y ayudantías.

\subsection{Distribución de las fechas}

Sofía Cisternas (DA) destaca que las solemnes quedaron injustamente distribuidas y que la mención Administración resultó con una peor distribución que la mención Economía. Por su parte, Mechones y Plan Común no presentan mayores quejas respecto de la distribución de las fechas y solo solicitan un mayor plazo de rendición de las solemnes.

Alan explica que las solemnes se revisaron con un representante por mención y generación, quienes analizaron la distribución de las fechas. Señala que la disponibilidad de salas y la reubicación de estudiantes constituyen el principal problema al momento de configurar dicha distribución. Agrega que la distribución fue evaluada mediante un programa que mide por pesos y realiza estimaciones y distribuciones más robustas según las distintas cantidades de participantes de cada ramo y su carga.

\subsection{Propuestas de las delegaturas}

\begin{itemize}
  \item \textbf{Participación de las menciones:} Sofía Cisternas propone que, en la configuración de las fechas, participe además un delegado de la mención correspondiente.
  \item \textbf{Mayor plazo de rendición:} Florencia Montiglio (DPC) solicita, a nombre de Mechones y Plan Común, un mayor plazo de rendición de las solemnes, alargando el tiempo de duración de los días de solemnes.
\end{itemize}

\subsection{Descoordinación entre evaluaciones, clases y ayudantías}

Se plantea que en TSI se produjo una diferencia sustancial de dificultad entre semestres, y se constata una descoordinación constante entre ayudantías, clases y solemnes. Esta situación se presenta en Finanzas I, TSI y Negocios II.

Asimismo, se señala que Tópicos Matemáticos, ramo de reforzamiento de matemáticas, va atrasado una semana respecto del ramo de Métodos Matemáticos. El Consejo discute que no tiene mucho sentido que el ramo de reforzamiento tenga contenidos atrasados respecto del ramo principal.

\subsection{Prerrequisito de TSI para Cuantitativo I}

Se expone que hay estudiantes a quienes no se les asignó cupo en TSI, ramo que es prerrequisito para Cuantitativo I, y que este último abre de todos modos. Ante esta situación, se consideran dos opciones:

\begin{enumerate}
  \item Quitar el prerrequisito de TSI a Cuantitativo I.
  \item Realizar un ``perdonazo'' a quienes reprobaron TSI y a quienes no tuvieron la oportunidad de tomar el ramo por falta de cupos.
\end{enumerate}

\subsection{Acuerdos y próximos pasos}

Ambas propuestas serán elevadas por la Secretaria de Asuntos Académicos, Manuela Fernández.

% ---------------------------------------------------------------------------
\section{Jaranazos}

\subsection{Antecedentes}

Se conversa sobre el Jaranazo de Agosto y los problemas asociados a este, y sobre cómo estos repercutieron en la planificación de las actividades masivas posteriores.

\subsection{FENDA}

Se recuerda que la FENDA fue un evento realizado en el contexto de las Fiestas Patrias en Chile, en el cual se montó una especie de fonda al interior de la FEN. El evento contó con varias actividades que resultaron con gran éxito, lo cual la comunidad hizo notar.

\subsection{Postura de la Escuela y acuerdo}

La Escuela quería prohibir la FENDA, debido a los problemas causados por el Jaranazo de Agosto y a la mala imagen que esta transmite de la institución. Frente a ello, los centros de estudiantes respondieron con una queja formal.

Finalmente se llegó al acuerdo de controlar el ingreso por medio de EARLY y de limitar el consumo a un máximo de dos cervezas por persona.

% ---------------------------------------------------------------------------
\section{Semana organizaciones}

Se informa que actualmente se está realizando la feria de organizaciones estudiantiles, instancia en la que las agrupaciones de la FEN dan a conocer su trabajo al estudiantado y que se enmarca en la Semana de Organizaciones.

% ---------------------------------------------------------------------------
\section{JOE}

Se informa con satisfacción que la FEN ganó los JOE, resultado que el Consejo celebra y reconoce como fruto de la participación de la comunidad.

% ---------------------------------------------------------------------------
\section{Reforma estatutos}

\subsection{Antecedentes}

Se retoma el proceso de reforma de los estatutos del CEIC, cuya propuesta se encuentra en elaboración.

\subsection{Acuerdos y próximos pasos}

\begin{enumerate}
  \item \textbf{Entrega de la propuesta:} La propuesta de reforma de estatutos se entregará con plazo límite el día 31 de octubre.
  \item \textbf{Consejo extraordinario:} Una vez entregada la propuesta, se solicitará un consejo extraordinario para votar el cambio de estatutos.
\end{enumerate}

% ---------------------------------------------------------------------------
\section{Varios}

\subsection{Denuncia ante el Tribunal de la FECH}

Se informa de una denuncia presentada ante el Tribunal de la FECH por el consejero FECH Nelson Guerra, quien señalaba que habían existido irregularidades en la forma en que se votaron los llamados a paro en el Consejo de Estudiantes de Ingeniería Comercial.

El Tribunal descartó las irregularidades y determinó que no hubo incumplimiento de los estatutos. Alan señala que todo este proceso no fue grato.

\vspace{0.5em}

Sin más puntos que tratar, se da por finalizada la sesión.

\vspace{4em}

\noindent\textbf{Firma:}

\vspace{3em}

\noindent Alan Evan Mc Coll Albuja\\
Presidente CEIC

\end{document}
"
%
% En la versión de revisión, todo lo marcado con \dato{...} aparece en color:
% son datos que deben verificarse o completarse antes de publicar el acta.

\documentclass[12pt,letterpaper]{article}

\usepackage[T1]{fontenc}
\usepackage[utf8]{inputenc}
\usepackage[spanish,es-nolists,es-nodecimaldot]{babel}
\usepackage[letterpaper,left=1in,right=1in,top=1.35in,bottom=1.1in,headheight=48pt,headsep=24pt,footskip=36pt]{geometry}
\usepackage{graphicx}
\usepackage{amssymb}
\usepackage{xcolor}
\usepackage{fancyhdr}
\usepackage{lastpage}
\usepackage{titlesec}
\usepackage{enumitem}
\usepackage{textcomp}

% ---------------------------------------------------------------------------
% Datos a verificar
% ---------------------------------------------------------------------------
\newif\ifrevision
\ifdefined\revision \revisiontrue \else \revisionfalse \fi
\newcommand{\dato}[1]{\ifrevision\textcolor[HTML]{C0392B}{#1}\else#1\fi}

% ---------------------------------------------------------------------------
% Datos del acta
% ---------------------------------------------------------------------------
\newcommand{\numeroacta}{8}
\newcommand{\fechaacta}{08-10-2026}

% ---------------------------------------------------------------------------
% Formato
% ---------------------------------------------------------------------------
\setlength{\parindent}{1.5em}
\setlength{\parskip}{6pt}

\pagestyle{fancy}
\fancyhf{}
\fancyhead[L]{\raisebox{-0.35\height}{\includegraphics[height=1.15cm]{logo_ceic.png}}}
\fancyhead[R]{Acta N\textdegree\numeroacta\\\fechaacta}
\fancyfoot[C]{Página \thepage\ de \pageref{LastPage}}
\renewcommand{\headrulewidth}{0.4pt}
\renewcommand{\footrulewidth}{0.4pt}

\titleformat{\section}{\normalfont\Large\bfseries}{\thesection.}{0.35em}{}
\titleformat{\subsection}{\normalfont\large\bfseries}{\thesubsection.}{0.35em}{}
\titlespacing*{\section}{0pt}{3.0ex plus 1ex minus .2ex}{1.6ex plus .2ex}
\titlespacing*{\subsection}{0pt}{2.4ex plus 1ex minus .2ex}{1.2ex plus .2ex}

\renewcommand{\labelitemi}{\tiny$\blacksquare$}
\setlist[itemize]{leftmargin=2.2em,itemsep=2pt}
\setlist[enumerate]{leftmargin=2.2em,itemsep=2pt}

\begin{document}

% ---------------------------------------------------------------------------
% Encabezado
% ---------------------------------------------------------------------------
\begin{center}
  {\Large\bfseries Centro de Estudiantes de Ingeniería Comercial}\\
  Universidad de Chile, Facultad de Economía y Negocios

  \vspace{1.2em}
  \textbf{Acta Consejo}
\end{center}

\vspace{0.8em}

\textbf{Asistentes:} Alan Mc Coll (Mesa), Diego Pakomio (Mesa extendida), Anabelee Araneda (DM), Arturo Mora (DM), Juan Budge (DM), Martín Farías (DPC), Florencia Montiglio (DPC), Sofía Cisternas (DA), Nicolás Sanzana (DE), Sofía Allende (CT).

\textbf{Hora de inicio:} 14:00 hrs

\textbf{Hora de término:} 15:20 hrs

\vspace{1em}

\textbf{Tabla:}

\begin{enumerate}[itemsep=6pt]
  \item Finanzas transparentes
  \item Solemnes
  \item Jaranazos
  \item Semana organizaciones
  \item JOE
  \item Reforma estatutos
  \item Varios
\end{enumerate}

% ---------------------------------------------------------------------------
\section{Finanzas transparentes}

\subsection{Balance presentado}

Se presenta al Consejo el balance correspondiente a los meses de Agosto y Septiembre, junto con el detalle de las principales fuentes de ingreso y categorías de gasto registradas en el libro diario de cada mes.

\begin{itemize}
  \item \textbf{Saldo del mes de Agosto (Ingresos $-$ Gastos):} \$161.503
  \item \textbf{Saldo del mes de Septiembre (Ingresos $-$ Gastos):} \$-682.811
\end{itemize}

Lo que arroja un saldo actual de caja chica de: \$9.820.820

\subsection{Detalle del periodo}

En Agosto, los ingresos ascendieron a \$261.840 y provinieron principalmente de las entradas de Cupido FEN (\$129.850), de la repartición de ganancias del Jaranazo (\$101.990) y de la participación del CEIC en la charla de diversidad realizada junto a Fundación Todo Mejora (\$30.000). Los gastos del mes alcanzaron los \$100.337 y correspondieron a la realización de Cupido FEN y del Jaranazo, a las libretas y stickers de la SUE para la feria de organizaciones y a gastos menores de funcionamiento.

En Septiembre, los ingresos fueron de \$221.550, compuestos por las entradas SIPEE (\$147.000) y por la devolución de una transferencia realizada por error (\$74.550). Los egresos sumaron \$904.361, entre los cuales destacan los gastos de la FENDA (guardias y DJ, entre otros), las dos jornadas de Tiny Desk, las libretas y stickers de la SUE, Cupido FEN y los desayunos realizados para deportistas y funcionarios.

\textbf{Nota:} Cinco de las entradas SIPEE se pagaron a \$7.000 en lugar de \$8.000 por un error de formulario; los \$1.000 restantes de cada una aparecerán en el balance de Octubre. Se recuerda que cualquier estudiante puede solicitar el detalle y los comprobantes de los balances presentados a través del formulario del Portal de Transparencia del CEIC.

% ---------------------------------------------------------------------------
\section{Solemnes}

\subsection{Antecedentes}

Se abre la discusión sobre la distribución de las fechas de las solemnes del semestre y su impacto en las distintas menciones y generaciones, así como sobre la coordinación general de las evaluaciones con las clases y ayudantías.

\subsection{Distribución de las fechas}

Sofía Cisternas (DA) destaca que las solemnes quedaron injustamente distribuidas y que la mención Administración resultó con una peor distribución que la mención Economía. Por su parte, Mechones y Plan Común no presentan mayores quejas respecto de la distribución de las fechas y solo solicitan un mayor plazo de rendición de las solemnes.

Alan explica que las solemnes se revisaron con un representante por mención y generación, quienes analizaron la distribución de las fechas. Señala que la disponibilidad de salas y la reubicación de estudiantes constituyen el principal problema al momento de configurar dicha distribución. Agrega que la distribución fue evaluada mediante un programa que mide por pesos y realiza estimaciones y distribuciones más robustas según las distintas cantidades de participantes de cada ramo y su carga.

\subsection{Propuestas de las delegaturas}

\begin{itemize}
  \item \textbf{Participación de las menciones:} Sofía Cisternas propone que, en la configuración de las fechas, participe además un delegado de la mención correspondiente.
  \item \textbf{Mayor plazo de rendición:} Florencia Montiglio (DPC) solicita, a nombre de Mechones y Plan Común, un mayor plazo de rendición de las solemnes, alargando el tiempo de duración de los días de solemnes.
\end{itemize}

\subsection{Descoordinación entre evaluaciones, clases y ayudantías}

Se plantea que en TSI se produjo una diferencia sustancial de dificultad entre semestres, y se constata una descoordinación constante entre ayudantías, clases y solemnes. Esta situación se presenta en Finanzas I, TSI y Negocios II.

Asimismo, se señala que Tópicos Matemáticos, ramo de reforzamiento de matemáticas, va atrasado una semana respecto del ramo de Métodos Matemáticos. El Consejo discute que no tiene mucho sentido que el ramo de reforzamiento tenga contenidos atrasados respecto del ramo principal.

\subsection{Prerrequisito de TSI para Cuantitativo I}

Se expone que hay estudiantes a quienes no se les asignó cupo en TSI, ramo que es prerrequisito para Cuantitativo I, y que este último abre de todos modos. Ante esta situación, se consideran dos opciones:

\begin{enumerate}
  \item Quitar el prerrequisito de TSI a Cuantitativo I.
  \item Realizar un ``perdonazo'' a quienes reprobaron TSI y a quienes no tuvieron la oportunidad de tomar el ramo por falta de cupos.
\end{enumerate}

\subsection{Acuerdos y próximos pasos}

Ambas propuestas serán elevadas por la Secretaria de Asuntos Académicos, Manuela Fernández.

% ---------------------------------------------------------------------------
\section{Jaranazos}

\subsection{Antecedentes}

Se conversa sobre el Jaranazo de Agosto y los problemas asociados a este, y sobre cómo estos repercutieron en la planificación de las actividades masivas posteriores.

\subsection{FENDA}

Se recuerda que la FENDA fue un evento realizado en el contexto de las Fiestas Patrias en Chile, en el cual se montó una especie de fonda al interior de la FEN. El evento contó con varias actividades que resultaron con gran éxito, lo cual la comunidad hizo notar.

\subsection{Postura de la Escuela y acuerdo}

La Escuela quería prohibir la FENDA, debido a los problemas causados por el Jaranazo de Agosto y a la mala imagen que esta transmite de la institución. Frente a ello, los centros de estudiantes respondieron con una queja formal.

Finalmente se llegó al acuerdo de controlar el ingreso por medio de EARLY y de limitar el consumo a un máximo de dos cervezas por persona.

% ---------------------------------------------------------------------------
\section{Semana organizaciones}

Se informa que actualmente se está realizando la feria de organizaciones estudiantiles, instancia en la que las agrupaciones de la FEN dan a conocer su trabajo al estudiantado y que se enmarca en la Semana de Organizaciones.

% ---------------------------------------------------------------------------
\section{JOE}

Se informa con satisfacción que la FEN ganó los JOE, resultado que el Consejo celebra y reconoce como fruto de la participación de la comunidad.

% ---------------------------------------------------------------------------
\section{Reforma estatutos}

\subsection{Antecedentes}

Se retoma el proceso de reforma de los estatutos del CEIC, cuya propuesta se encuentra en elaboración.

\subsection{Acuerdos y próximos pasos}

\begin{enumerate}
  \item \textbf{Entrega de la propuesta:} La propuesta de reforma de estatutos se entregará con plazo límite el día 31 de octubre.
  \item \textbf{Consejo extraordinario:} Una vez entregada la propuesta, se solicitará un consejo extraordinario para votar el cambio de estatutos.
\end{enumerate}

% ---------------------------------------------------------------------------
\section{Varios}

\subsection{Denuncia ante el Tribunal de la FECH}

Se informa de una denuncia presentada ante el Tribunal de la FECH por el consejero FECH Nelson Guerra, quien señalaba que habían existido irregularidades en la forma en que se votaron los llamados a paro en el Consejo de Estudiantes de Ingeniería Comercial.

El Tribunal descartó las irregularidades y determinó que no hubo incumplimiento de los estatutos. Alan señala que todo este proceso no fue grato.

\vspace{0.5em}

Sin más puntos que tratar, se da por finalizada la sesión.

\vspace{4em}

\noindent\textbf{Firma:}

\vspace{3em}

\noindent Alan Evan Mc Coll Albuja\\
Presidente CEIC

\end{document}
"
%
% En la versión de revisión, todo lo marcado con \dato{...} aparece en color:
% son datos que deben verificarse o completarse antes de publicar el acta.

\documentclass[12pt,letterpaper]{article}

\usepackage[T1]{fontenc}
\usepackage[utf8]{inputenc}
\usepackage[spanish,es-nolists,es-nodecimaldot]{babel}
\usepackage[letterpaper,left=1in,right=1in,top=1.35in,bottom=1.1in,headheight=48pt,headsep=24pt,footskip=36pt]{geometry}
\usepackage{graphicx}
\usepackage{amssymb}
\usepackage{xcolor}
\usepackage{fancyhdr}
\usepackage{lastpage}
\usepackage{titlesec}
\usepackage{enumitem}
\usepackage{textcomp}

% ---------------------------------------------------------------------------
% Datos a verificar
% ---------------------------------------------------------------------------
\newif\ifrevision
\ifdefined\revision \revisiontrue \else \revisionfalse \fi
\newcommand{\dato}[1]{\ifrevision\textcolor[HTML]{C0392B}{#1}\else#1\fi}

% ---------------------------------------------------------------------------
% Datos del acta
% ---------------------------------------------------------------------------
\newcommand{\numeroacta}{8}
\newcommand{\fechaacta}{08-10-2026}

% ---------------------------------------------------------------------------
% Formato
% ---------------------------------------------------------------------------
\setlength{\parindent}{1.5em}
\setlength{\parskip}{6pt}

\pagestyle{fancy}
\fancyhf{}
\fancyhead[L]{\raisebox{-0.35\height}{\includegraphics[height=1.15cm]{logo_ceic.png}}}
\fancyhead[R]{Acta N\textdegree\numeroacta\\\fechaacta}
\fancyfoot[C]{Página \thepage\ de \pageref{LastPage}}
\renewcommand{\headrulewidth}{0.4pt}
\renewcommand{\footrulewidth}{0.4pt}

\titleformat{\section}{\normalfont\Large\bfseries}{\thesection.}{0.35em}{}
\titleformat{\subsection}{\normalfont\large\bfseries}{\thesubsection.}{0.35em}{}
\titlespacing*{\section}{0pt}{3.0ex plus 1ex minus .2ex}{1.6ex plus .2ex}
\titlespacing*{\subsection}{0pt}{2.4ex plus 1ex minus .2ex}{1.2ex plus .2ex}

\renewcommand{\labelitemi}{\tiny$\blacksquare$}
\setlist[itemize]{leftmargin=2.2em,itemsep=2pt}
\setlist[enumerate]{leftmargin=2.2em,itemsep=2pt}

\begin{document}

% ---------------------------------------------------------------------------
% Encabezado
% ---------------------------------------------------------------------------
\begin{center}
  {\Large\bfseries Centro de Estudiantes de Ingeniería Comercial}\\
  Universidad de Chile, Facultad de Economía y Negocios

  \vspace{1.2em}
  \textbf{Acta Consejo}
\end{center}

\vspace{0.8em}

\textbf{Asistentes:} Alan Mc Coll (Mesa), Diego Pakomio (Mesa extendida), Anabelee Araneda (DM), Arturo Mora (DM), Juan Budge (DM), Martín Farías (DPC), Florencia Montiglio (DPC), Sofía Cisternas (DA), Nicolás Sanzana (DE), Sofía Allende (CT).

\textbf{Hora de inicio:} 14:00 hrs

\textbf{Hora de término:} 15:20 hrs

\vspace{1em}

\textbf{Tabla:}

\begin{enumerate}[itemsep=6pt]
  \item Finanzas transparentes
  \item Solemnes
  \item Jaranazos
  \item Semana organizaciones
  \item JOE
  \item Reforma estatutos
  \item Varios
\end{enumerate}

% ---------------------------------------------------------------------------
\section{Finanzas transparentes}

\subsection{Balance presentado}

Se presenta al Consejo el balance correspondiente a los meses de Agosto y Septiembre, junto con el detalle de las principales fuentes de ingreso y categorías de gasto registradas en el libro diario de cada mes.

\begin{itemize}
  \item \textbf{Saldo del mes de Agosto (Ingresos $-$ Gastos):} \$161.503
  \item \textbf{Saldo del mes de Septiembre (Ingresos $-$ Gastos):} \$-682.811
\end{itemize}

Lo que arroja un saldo actual de caja chica de: \$9.820.820

\subsection{Detalle del periodo}

En Agosto, los ingresos ascendieron a \$261.840 y provinieron principalmente de las entradas de Cupido FEN (\$129.850), de la repartición de ganancias del Jaranazo (\$101.990) y de la participación del CEIC en la charla de diversidad realizada junto a Fundación Todo Mejora (\$30.000). Los gastos del mes alcanzaron los \$100.337 y correspondieron a la realización de Cupido FEN y del Jaranazo, a las libretas y stickers de la SUE para la feria de organizaciones y a gastos menores de funcionamiento.

En Septiembre, los ingresos fueron de \$221.550, compuestos por las entradas SIPEE (\$147.000) y por la devolución de una transferencia realizada por error (\$74.550). Los egresos sumaron \$904.361, entre los cuales destacan los gastos de la FENDA (guardias y DJ, entre otros), las dos jornadas de Tiny Desk, las libretas y stickers de la SUE, Cupido FEN y los desayunos realizados para deportistas y funcionarios.

\textbf{Nota:} Cinco de las entradas SIPEE se pagaron a \$7.000 en lugar de \$8.000 por un error de formulario; los \$1.000 restantes de cada una aparecerán en el balance de Octubre. Se recuerda que cualquier estudiante puede solicitar el detalle y los comprobantes de los balances presentados a través del formulario del Portal de Transparencia del CEIC.

% ---------------------------------------------------------------------------
\section{Solemnes}

\subsection{Antecedentes}

Se abre la discusión sobre la distribución de las fechas de las solemnes del semestre y su impacto en las distintas menciones y generaciones, así como sobre la coordinación general de las evaluaciones con las clases y ayudantías.

\subsection{Distribución de las fechas}

Sofía Cisternas (DA) destaca que las solemnes quedaron injustamente distribuidas y que la mención Administración resultó con una peor distribución que la mención Economía. Por su parte, Mechones y Plan Común no presentan mayores quejas respecto de la distribución de las fechas y solo solicitan un mayor plazo de rendición de las solemnes.

Alan explica que las solemnes se revisaron con un representante por mención y generación, quienes analizaron la distribución de las fechas. Señala que la disponibilidad de salas y la reubicación de estudiantes constituyen el principal problema al momento de configurar dicha distribución. Agrega que la distribución fue evaluada mediante un programa que mide por pesos y realiza estimaciones y distribuciones más robustas según las distintas cantidades de participantes de cada ramo y su carga.

\subsection{Propuestas de las delegaturas}

\begin{itemize}
  \item \textbf{Participación de las menciones:} Sofía Cisternas propone que, en la configuración de las fechas, participe además un delegado de la mención correspondiente.
  \item \textbf{Mayor plazo de rendición:} Florencia Montiglio (DPC) solicita, a nombre de Mechones y Plan Común, un mayor plazo de rendición de las solemnes, alargando el tiempo de duración de los días de solemnes.
\end{itemize}

\subsection{Descoordinación entre evaluaciones, clases y ayudantías}

Se plantea que en TSI se produjo una diferencia sustancial de dificultad entre semestres, y se constata una descoordinación constante entre ayudantías, clases y solemnes. Esta situación se presenta en Finanzas I, TSI y Negocios II.

Asimismo, se señala que Tópicos Matemáticos, ramo de reforzamiento de matemáticas, va atrasado una semana respecto del ramo de Métodos Matemáticos. El Consejo discute que no tiene mucho sentido que el ramo de reforzamiento tenga contenidos atrasados respecto del ramo principal.

\subsection{Prerrequisito de TSI para Cuantitativo I}

Se expone que hay estudiantes a quienes no se les asignó cupo en TSI, ramo que es prerrequisito para Cuantitativo I, y que este último abre de todos modos. Ante esta situación, se consideran dos opciones:

\begin{enumerate}
  \item Quitar el prerrequisito de TSI a Cuantitativo I.
  \item Realizar un ``perdonazo'' a quienes reprobaron TSI y a quienes no tuvieron la oportunidad de tomar el ramo por falta de cupos.
\end{enumerate}

\subsection{Acuerdos y próximos pasos}

Ambas propuestas serán elevadas por la Secretaria de Asuntos Académicos, Manuela Fernández.

% ---------------------------------------------------------------------------
\section{Jaranazos}

\subsection{Antecedentes}

Se conversa sobre el Jaranazo de Agosto y los problemas asociados a este, y sobre cómo estos repercutieron en la planificación de las actividades masivas posteriores.

\subsection{FENDA}

Se recuerda que la FENDA fue un evento realizado en el contexto de las Fiestas Patrias en Chile, en el cual se montó una especie de fonda al interior de la FEN. El evento contó con varias actividades que resultaron con gran éxito, lo cual la comunidad hizo notar.

\subsection{Postura de la Escuela y acuerdo}

La Escuela quería prohibir la FENDA, debido a los problemas causados por el Jaranazo de Agosto y a la mala imagen que esta transmite de la institución. Frente a ello, los centros de estudiantes respondieron con una queja formal.

Finalmente se llegó al acuerdo de controlar el ingreso por medio de EARLY y de limitar el consumo a un máximo de dos cervezas por persona.

% ---------------------------------------------------------------------------
\section{Semana organizaciones}

Se informa que actualmente se está realizando la feria de organizaciones estudiantiles, instancia en la que las agrupaciones de la FEN dan a conocer su trabajo al estudiantado y que se enmarca en la Semana de Organizaciones.

% ---------------------------------------------------------------------------
\section{JOE}

Se informa con satisfacción que la FEN ganó los JOE, resultado que el Consejo celebra y reconoce como fruto de la participación de la comunidad.

% ---------------------------------------------------------------------------
\section{Reforma estatutos}

\subsection{Antecedentes}

Se retoma el proceso de reforma de los estatutos del CEIC, cuya propuesta se encuentra en elaboración.

\subsection{Acuerdos y próximos pasos}

\begin{enumerate}
  \item \textbf{Entrega de la propuesta:} La propuesta de reforma de estatutos se entregará con plazo límite el día 31 de octubre.
  \item \textbf{Consejo extraordinario:} Una vez entregada la propuesta, se solicitará un consejo extraordinario para votar el cambio de estatutos.
\end{enumerate}

% ---------------------------------------------------------------------------
\section{Varios}

\subsection{Denuncia ante el Tribunal de la FECH}

Se informa de una denuncia presentada ante el Tribunal de la FECH por el consejero FECH Nelson Guerra, quien señalaba que habían existido irregularidades en la forma en que se votaron los llamados a paro en el Consejo de Estudiantes de Ingeniería Comercial.

El Tribunal descartó las irregularidades y determinó que no hubo incumplimiento de los estatutos. Alan señala que todo este proceso no fue grato.

\vspace{0.5em}

Sin más puntos que tratar, se da por finalizada la sesión.

\vspace{4em}

\noindent\textbf{Firma:}

\vspace{3em}

\noindent Alan Evan Mc Coll Albuja\\
Presidente CEIC

\end{document}
"
%
% En la versión de revisión, todo lo marcado con \dato{...} aparece en color:
% son datos que deben verificarse o completarse antes de publicar el acta.

\documentclass[12pt,letterpaper]{article}

\usepackage[T1]{fontenc}
\usepackage[utf8]{inputenc}
\usepackage[spanish,es-nolists,es-nodecimaldot]{babel}
\usepackage[letterpaper,left=1in,right=1in,top=1.35in,bottom=1.1in,headheight=48pt,headsep=24pt,footskip=36pt]{geometry}
\usepackage{graphicx}
\usepackage{amssymb}
\usepackage{xcolor}
\usepackage{fancyhdr}
\usepackage{lastpage}
\usepackage{titlesec}
\usepackage{enumitem}
\usepackage{textcomp}

% ---------------------------------------------------------------------------
% Datos a verificar
% ---------------------------------------------------------------------------
\newif\ifrevision
\ifdefined\revision \revisiontrue \else \revisionfalse \fi
\newcommand{\dato}[1]{\ifrevision\textcolor[HTML]{C0392B}{#1}\else#1\fi}

% ---------------------------------------------------------------------------
% Datos del acta
% ---------------------------------------------------------------------------
\newcommand{\numeroacta}{8}
\newcommand{\fechaacta}{08-10-2026}

% ---------------------------------------------------------------------------
% Formato
% ---------------------------------------------------------------------------
\setlength{\parindent}{1.5em}
\setlength{\parskip}{6pt}

\pagestyle{fancy}
\fancyhf{}
\fancyhead[L]{\raisebox{-0.35\height}{\includegraphics[height=1.15cm]{logo_ceic.png}}}
\fancyhead[R]{Acta N\textdegree\numeroacta\\\fechaacta}
\fancyfoot[C]{Página \thepage\ de \pageref{LastPage}}
\renewcommand{\headrulewidth}{0.4pt}
\renewcommand{\footrulewidth}{0.4pt}

\titleformat{\section}{\normalfont\Large\bfseries}{\thesection.}{0.35em}{}
\titleformat{\subsection}{\normalfont\large\bfseries}{\thesubsection.}{0.35em}{}
\titlespacing*{\section}{0pt}{3.0ex plus 1ex minus .2ex}{1.6ex plus .2ex}
\titlespacing*{\subsection}{0pt}{2.4ex plus 1ex minus .2ex}{1.2ex plus .2ex}

\renewcommand{\labelitemi}{\tiny$\blacksquare$}
\setlist[itemize]{leftmargin=2.2em,itemsep=2pt}
\setlist[enumerate]{leftmargin=2.2em,itemsep=2pt}

\begin{document}

% ---------------------------------------------------------------------------
% Encabezado
% ---------------------------------------------------------------------------
\begin{center}
  {\Large\bfseries Centro de Estudiantes de Ingeniería Comercial}\\
  Universidad de Chile, Facultad de Economía y Negocios

  \vspace{1.2em}
  \textbf{Acta Consejo}
\end{center}

\vspace{0.8em}

\textbf{Asistentes:} Alan Mc Coll (Mesa), Diego Pakomio (Mesa extendida), Anabelee Araneda (DM), Arturo Mora (DM), Juan Budge (DM), Martín Farías (DPC), Florencia Montiglio (DPC), Sofía Cisternas (DA), Nicolás Sanzana (DE), Sofía Allende (CT).

\textbf{Hora de inicio:} 14:00 hrs

\textbf{Hora de término:} 15:20 hrs

\vspace{1em}

\textbf{Tabla:}

\begin{enumerate}[itemsep=6pt]
  \item Finanzas transparentes
  \item Solemnes
  \item Jaranazos
  \item Semana organizaciones
  \item JOE
  \item Reforma estatutos
  \item Varios
\end{enumerate}

% ---------------------------------------------------------------------------
\section{Finanzas transparentes}

\subsection{Balance presentado}

Se presenta al Consejo el balance correspondiente a los meses de Agosto y Septiembre, junto con el detalle de las principales fuentes de ingreso y categorías de gasto registradas en el libro diario de cada mes.

\begin{itemize}
  \item \textbf{Saldo del mes de Agosto (Ingresos $-$ Gastos):} \$161.503
  \item \textbf{Saldo del mes de Septiembre (Ingresos $-$ Gastos):} \$-682.811
\end{itemize}

Lo que arroja un saldo actual de caja chica de: \$9.820.820

\subsection{Detalle del periodo}

En Agosto, los ingresos ascendieron a \$261.840 y provinieron principalmente de las entradas de Cupido FEN (\$129.850), de la repartición de ganancias del Jaranazo (\$101.990) y de la participación del CEIC en la charla de diversidad realizada junto a Fundación Todo Mejora (\$30.000). Los gastos del mes alcanzaron los \$100.337 y correspondieron a la realización de Cupido FEN y del Jaranazo, a las libretas y stickers de la SUE para la feria de organizaciones y a gastos menores de funcionamiento.

En Septiembre, los ingresos fueron de \$221.550, compuestos por las entradas SIPEE (\$147.000) y por la devolución de una transferencia realizada por error (\$74.550). Los egresos sumaron \$904.361, entre los cuales destacan los gastos de la FENDA (guardias y DJ, entre otros), las dos jornadas de Tiny Desk, las libretas y stickers de la SUE, Cupido FEN y los desayunos realizados para deportistas y funcionarios.

\textbf{Nota:} Cinco de las entradas SIPEE se pagaron a \$7.000 en lugar de \$8.000 por un error de formulario; los \$1.000 restantes de cada una aparecerán en el balance de Octubre. Se recuerda que cualquier estudiante puede solicitar el detalle y los comprobantes de los balances presentados a través del formulario del Portal de Transparencia del CEIC.

% ---------------------------------------------------------------------------
\section{Solemnes}

\subsection{Antecedentes}

Se abre la discusión sobre la distribución de las fechas de las solemnes del semestre y su impacto en las distintas menciones y generaciones, así como sobre la coordinación general de las evaluaciones con las clases y ayudantías.

\subsection{Distribución de las fechas}

Sofía Cisternas (DA) destaca que las solemnes quedaron injustamente distribuidas y que la mención Administración resultó con una peor distribución que la mención Economía. Por su parte, Mechones y Plan Común no presentan mayores quejas respecto de la distribución de las fechas y solo solicitan un mayor plazo de rendición de las solemnes.

Alan explica que las solemnes se revisaron con un representante por mención y generación, quienes analizaron la distribución de las fechas. Señala que la disponibilidad de salas y la reubicación de estudiantes constituyen el principal problema al momento de configurar dicha distribución. Agrega que la distribución fue evaluada mediante un programa que mide por pesos y realiza estimaciones y distribuciones más robustas según las distintas cantidades de participantes de cada ramo y su carga.

\subsection{Propuestas de las delegaturas}

\begin{itemize}
  \item \textbf{Participación de las menciones:} Sofía Cisternas propone que, en la configuración de las fechas, participe además un delegado de la mención correspondiente.
  \item \textbf{Mayor plazo de rendición:} Florencia Montiglio (DPC) solicita, a nombre de Mechones y Plan Común, un mayor plazo de rendición de las solemnes, alargando el tiempo de duración de los días de solemnes.
\end{itemize}

\subsection{Descoordinación entre evaluaciones, clases y ayudantías}

Se plantea que en TSI se produjo una diferencia sustancial de dificultad entre semestres, y se constata una descoordinación constante entre ayudantías, clases y solemnes. Esta situación se presenta en Finanzas I, TSI y Negocios II.

Asimismo, se señala que Tópicos Matemáticos, ramo de reforzamiento de matemáticas, va atrasado una semana respecto del ramo de Métodos Matemáticos. El Consejo discute que no tiene mucho sentido que el ramo de reforzamiento tenga contenidos atrasados respecto del ramo principal.

\subsection{Prerrequisito de TSI para Cuantitativo I}

Se expone que hay estudiantes a quienes no se les asignó cupo en TSI, ramo que es prerrequisito para Cuantitativo I, y que este último abre de todos modos. Ante esta situación, se consideran dos opciones:

\begin{enumerate}
  \item Quitar el prerrequisito de TSI a Cuantitativo I.
  \item Realizar un ``perdonazo'' a quienes reprobaron TSI y a quienes no tuvieron la oportunidad de tomar el ramo por falta de cupos.
\end{enumerate}

\subsection{Acuerdos y próximos pasos}

Ambas propuestas serán elevadas por la Secretaria de Asuntos Académicos, Manuela Fernández.

% ---------------------------------------------------------------------------
\section{Jaranazos}

\subsection{Antecedentes}

Se conversa sobre el Jaranazo de Agosto y los problemas asociados a este, y sobre cómo estos repercutieron en la planificación de las actividades masivas posteriores.

\subsection{FENDA}

Se recuerda que la FENDA fue un evento realizado en el contexto de las Fiestas Patrias en Chile, en el cual se montó una especie de fonda al interior de la FEN. El evento contó con varias actividades que resultaron con gran éxito, lo cual la comunidad hizo notar.

\subsection{Postura de la Escuela y acuerdo}

La Escuela quería prohibir la FENDA, debido a los problemas causados por el Jaranazo de Agosto y a la mala imagen que esta transmite de la institución. Frente a ello, los centros de estudiantes respondieron con una queja formal.

Finalmente se llegó al acuerdo de controlar el ingreso por medio de EARLY y de limitar el consumo a un máximo de dos cervezas por persona.

% ---------------------------------------------------------------------------
\section{Semana organizaciones}

Se informa que actualmente se está realizando la feria de organizaciones estudiantiles, instancia en la que las agrupaciones de la FEN dan a conocer su trabajo al estudiantado y que se enmarca en la Semana de Organizaciones.

% ---------------------------------------------------------------------------
\section{JOE}

Se informa con satisfacción que la FEN ganó los JOE, resultado que el Consejo celebra y reconoce como fruto de la participación de la comunidad.

% ---------------------------------------------------------------------------
\section{Reforma estatutos}

\subsection{Antecedentes}

Se retoma el proceso de reforma de los estatutos del CEIC, cuya propuesta se encuentra en elaboración.

\subsection{Acuerdos y próximos pasos}

\begin{enumerate}
  \item \textbf{Entrega de la propuesta:} La propuesta de reforma de estatutos se entregará con plazo límite el día 31 de octubre.
  \item \textbf{Consejo extraordinario:} Una vez entregada la propuesta, se solicitará un consejo extraordinario para votar el cambio de estatutos.
\end{enumerate}

% ---------------------------------------------------------------------------
\section{Varios}

\subsection{Denuncia ante el Tribunal de la FECH}

Se informa de una denuncia presentada ante el Tribunal de la FECH por el consejero FECH Nelson Guerra, quien señalaba que habían existido irregularidades en la forma en que se votaron los llamados a paro en el Consejo de Estudiantes de Ingeniería Comercial.

El Tribunal descartó las irregularidades y determinó que no hubo incumplimiento de los estatutos. Alan señala que todo este proceso no fue grato.

\vspace{0.5em}

Sin más puntos que tratar, se da por finalizada la sesión.

\vspace{4em}

\noindent\textbf{Firma:}

\vspace{3em}

\noindent Alan Evan Mc Coll Albuja\\
Presidente CEIC

\end{document}
